\documentclass[slides]{pgnotes}

\title{Time Series}

\begin{document}

\maketitle

\section{Time series data}

A time series consists of data points indexed in time order.
Encountered in many contexts.
Examples: 

\begin{enumerate}
\item Stock price of a given stock or consolidated index at the close of business each day.
\item Number of entries and exits recorded by a car park barrier system as logged at 11PM each night.
\item Air temperature, pressure, UV and humidity recorded every 15 minutes at 20 unmanned weather stations. 
\item Number of cancelled trains as a percentage of scheduled services on a rail line between Dublin and Belfast per year since 1920. 
\end{enumerate}

\subsection{Components}

\begin{description}
\item[Trend:] a systematic change following a linear / nonlinear path, changing over time, aperiodic.
\item[Seasonality:] a systematic change following a linear / nonlinear path, changing over time, periodic
\item[Cyclic:] systematic change that repeats though not fixed / indexed to time. 
\item[Noise:] non-systematic component (i.e. not a trend or seasonal variation) within the data
\end{description}

\subsection{Analysis tools}

A number of methods are available for time series analysis, including forecasting likely future behaviour.
Practical implementations are found in R, Pandas, MATLAB and other packages.
However, it is also possible to construct these analyses within the database layer. 

This class does not focus on the theory of these analysis tools.
Rather their application / implementation within the database (as opposed to at the application layer) is discussed.

\subsection{Sample Dataset}

The following demonstration of time series analysis will use the \href{https://data.gov.ie/dataset/covidstatisticsprofilehpscirelandopendata1}{CovidStatisticsProfileHPSCIrelandOpenData} from \url{https://data.gov.ie}.


\section{Window functions}

We have already met aggregate functions: sum, min, max, others. 
These normally reduce a number of rows to a single row, or perhaps more if \texttt{GROUP BY} is present.

Time series analyses depend usually on rows relative to the current row. 
Examples:
\begin{itemize}
\item A cumalative sum of a particular expression
\item A moving average of a particular expression
\item The row number in the result set
\item The rank order of a particular row with respect to a particular expression. Includes numerical rank, percentage rank, percentage/normalised ranks, percentiles.
\item Value at a row offset before / after the current row.
\item First / last value in the window frame. 
\end{itemize}

\subsection{Window functions vs aggregate functions}

\href{https://www.postgresql.org/docs/13/tutorial-window.html}{Window functions} are like aggregate functions except:
\begin{itemize}
\item The aggregation occurs over a defined \textbf{window frame}, or set of rows relative to the current row, rather than over the whole table.
\item 
  Generally all aggregate functions can be used over a window, but there are some additional aggregate functions that only make sense over a window.
  \end{itemize}

\subsection{OVER clause}

A window function can be recognised by the \mintinline{postgresql}{OVER} clause, which may contain subclauses:
\begin{description}
\item[ORDER BY] to  (not mandatory but usually required)
\item[Frame clause] to define the window frame relative to the start / end.  If unspecified, it's by default from start to the current row.
\item[PARTITION BY] to break the analysis down by another column / expression
\end{description}

\subsection{OVER clause example}

Consider a situation where we wish to compute a cumalative sum with respect to each date.
In the sample dataset, we wish to provide a running total of the \texttt{confirmed\_cases} column.
The following SQL would realise this: 

\begin{minted}{postgresql}
SELECT report_date, confirmed_cases,
SUM(confirmed_cases) OVER ( ORDER BY report_date ASC ) as total_confirmed_cases
ORDER BY report_date asc ;
\end{minted}
Points to note:
\begin{itemize}
\item
  We want to use the \texttt{SUM} aggregate function \texttt{OVER} the case total to date.
  This means that we have an \texttt{ORDER BY} within the \texttt{OVER} clause.
\item 
  The main \texttt{ORDER BY} clause is independent of that in the \texttt{OVER} clause.
  Although often the same as that within the \texttt{OVER} clause(s) it can differ.
\end{itemize}

\subsection{Frame specification}

Consider now that we want the moving average.
This is similar to the cumulative sum, except we must define a particular window over which the average is computed.
For example, a 6-point moving average: 
\begin{minted}{postgresql}
SELECT report_date, confirmed_cases, 
CAST(
AVG(confirmed_cases) 
OVER (ORDER BY report_date ASC ROWS BETWEEN 6 PRECEDING AND CURRENT ROW ) 
as integer ) as ma7d
FROM covid
ORDER BY report_date ASC ;
\end{minted}

\newpage
Points to note: 
\begin{itemize}
\item The \textbf{frame specification} \mintinline{postgresql}{ROWS BETWEEN 6 PRECEDING AND CURRENT ROW}  defines the window frame, or set of rows related to the current row, to be used for the window / aggregate function. 
\item We use \mintinline{postgresql}{CAST} to coerce the floating point average back to an integer. May / may not be required depending on situation of interest.
\end{itemize}

\subsection{Window re-use}

We may wish to compute multiple metrics over the same window.
For example: 
\begin{minted}{postgresql}
SELECT report_date, confirmed_cases,
SUM(confirmed_cases) OVER ( ORDER BY report_date ASC ) as total_confirmed_cases,
SUM(confirmed_deaths) OVER ( ORDER BY report_date ASC ) as total_confirmed_deaths 
FROM covid 
ORDER BY report_date asc ;
\end{minted}
Notice here that window is the same.

\subsection{Named window}

Instead of defining it each time we can created a named window. 
\begin{minted}{postgresql}
SELECT report_date, confirmed_cases,
SUM(confirmed_cases) OVER reports_to_date as total_confirmed_cases,
SUM(confirmed_deaths) OVER reports_to_date as total_confirmed_deaths 
FROM covid 
WINDOW reports_to_date AS ( ORDER BY report_date ASC ) 
ORDER BY report_date asc ;
\end{minted}

\subsection{Multiple named window frames}

Multiple named window frames can be defined and used within the same query.
Pay attention to syntax for multiple windows: 
\begin{minted}{postgresql}
SELECT report_date, confirmed_cases,
CAST( AVG(confirmed_cases) OVER last_7_days as integer ) as cases_7d,
CAST( AVG(confirmed_cases) OVER last_14_days as integer ) as cases_14d,
CAST( AVG(confirmed_deaths) OVER last_7_days as integer ) as deaths_7d,
CAST( AVG(confirmed_deaths) OVER last_14_days as integer ) as deaths_14d
FROM covid 
WINDOW
last_7_days AS (ORDER BY report_date ASC ROWS BETWEEN 6 PRECEDING AND CURRENT ROW ), 
last_14_days AS (ORDER BY report_date ASC ROWS BETWEEN 13 PRECEDING AND CURRENT ROW ) 
ORDER BY report_date ASC ;
\end{minted} 




\end{document}
